\chapter{مفاهیم پایه‌ای}
\label{chap:concepts}

هدف این فصل معرفی مبانی نظری و فناوری‌های مورد استفاده در چارچوب و برنامه‌ی نمونه است. نخست سطوح مختلف تست نرم‌افزار و جایگاه تست سرتاسری روشن می‌شود؛ سپس سازوکار خودکارسازی رابط کاربری در ویندوز، الگوی شیء صفحه و زبان مخصوص دامنه، و در پایان فناوری‌های به‌کاررفته در برنامه‌ی نمونه و ابزارهای گزارش‌گیری بررسی می‌شوند.

\section{سطوح تست نرم‌افزار}
تست نرم‌افزار معمولاً در چند سطح سازمان می‌یابد: تست واحد که رفتار کوچک‌ترین اجزا را بررسی می‌کند؛ تست یکپارچگی که درستی تعامل اجزا با یکدیگر را می‌سنجد و تست سرتاسری
\lr{\gls{E2E}}
که سامانه را از دید کاربر نهایی و به صورت یکپارچه تست می‌کند
\cite{crispin2009}. هرچه از سطح واحد به سطح سرتاسری نزدیک‌تر می‌شویم، شبیه‌سازی واقعیت بیشتر می‌شود اما هزینه‌ی اجرا و نگهداری تست‌ها و شکنندگی آن‌ها در برابر تغییرات نیز افزایش می‌یابد. تست‌های سرتاسری کلاسیک به همین دلیل شکننده شناخته می‌شوند: کوچک‌ترین تغییر ظاهری رابط کاربری می‌تواند ده‌ها تست را از کار بیندازد. اصل طراحی چارچوب این پایان‌نامه، کاهش این شکنندگی از طریق قراردادهای پایداری و جداسازی لایه‌ها است که در فصل
\autoref{chap:framework}
تشریح می‌شود.

\section{خودکارسازی رابط کاربری در ویندوز}
\lr{\gls{UIA}}
رابط برنامه‌نویسی دسترس‌پذیری ویندوز است که درخت سلسله‌مراتبی عناصر رابط کاربری هر برنامه را در اختیار می‌گذارد
\cite{uia}. هر عنصر ویژگی‌هایی مانند نام
\lr{(Name)}
و شناسه‌ی خودکارسازی
\lr{(AutomationId)}
دارد و الگوهای تعامل مانند فراخوانی دکمه یا تنظیم مقدار را ارائه می‌کند. برنامه‌های
\lr{Qt}
این درخت را از روی ویژگی‌های
\lr{objectName}
و
\lr{accessibleName}
ساخت می‌دهند؛ این نکته پایه‌ی قرارداد «شناسه پایدار» در فصل
\autoref{chap:framework}
است.

\lr{WinAppDriver}
\cite{winappdriver}
درایور خودکارسازی مایکروسافت است که پروتکل
\lr{W3C WebDriver}
را برای برنامه‌های ویندوز پیاده می‌کند و از طریق
\lr{\gls{UIA}}
با عناصر رابط کاربری کار می‌کند.
\lr{Appium}
\cite{appium}
سرور پروتکل است که تست‌ها را از هر زبان برنامه‌نویسی به درایور مناسب ترجمه می‌کند. این ترکیب امکان می‌دهد تست‌ها با
\lr{Python}
نوشته شوند و برنامه‌ی دسکتاپ دقیقاً مانند یک برنامه‌ی موبایل یا وب خودکار شود.

\section{الگوی شیء صفحه و زبان مخصوص دامنه}
در الگوی شیء صفحه
\lr{\gls{POM}}
\cite{fowler2002}
برای هر صفحه یا پنجره‌ی برنامه، یک کلاس مسئول ساخته می‌شود که تنها نهادِ آن صفحه را می‌شناسد و تعاملات را پنهان می‌کند. تست‌ها با شیء صفحه گفت‌وگو می‌کنند و از جزئیات مکان‌یابی عناصر بی‌خبرند. لایه‌ی بالاتر، زبان مخصوص دامنه
\lr{\gls{DSL}}
است: مجموعه‌ای از متدهای معنادار برای حوزه‌ی مسئله که خواندن تست‌ها را به روایت رفتار کاربر تبدیل می‌کند. ترکیب این دو لایه یک اصل ساده می‌سازد: «فقط لایه‌ی پایین با درایور حرف می‌زند». به این ترتیب تست‌ها مستند زنده‌ی رفتار برنامه‌اند و نگهداری آن‌ها محدود به یک لایه است.

\section{فناوری‌های برنامه‌ی نمونه}
برنامه‌ی نمونه با
\lr{Qt Widgets}
نسخه‌ی ۶.۸ ساخته شده است
\cite{qtdocs};
کتابخانه‌ی استاندارد رابط کاربری دسکتاپ
\lr{Qt}
که درخت عناصر خود را برای
\lr{\gls{UIA}}
آشکار می‌کند. نمودارها با ماژول
\lr{Qt Charts}
ترسیم می‌شوند و داده‌ها در پایگاه‌داده‌ی
\lr{SQLite}
ذخیره می‌شوند
\cite{sqlite}.
لایه‌ی دسترسی به داده‌ی برنامه به صورت یک کلاس واحد است که منبع یگانه‌ی حقیقت دامنه محسوب می‌شود و هم رابط کاربری و هم تست‌های یکپارچگی از آن تغذیه می‌کنند. علاوه بر تست‌های سرتاسری، برنامه‌ی نمونه دو تست داخلی دارد: تست دود رابط کاربری و تست لایه‌ی داده که با
\lr{Qt Test}
\cite{qtcitest}
نوشته شده‌اند و در پایپلاین پیش از تست‌های سرتاسری اجرا می‌شوند.

\section{\lr{pytest} و گزارش‌گیری \lr{Allure}}
تست‌های چارچوب با فریم‌ورک
\lr{pytest}
\cite{pytest}
نوشته می‌شوند که از فیکسچرها
\gls{Fixture}
برای آماده‌سازی محیط هر تست بهره می‌گیرد. گزارش‌گیری با
\lr{Allure}
\cite{allure}
انجام می‌شود؛ هر گام تست، فایل پیوست (از جمله تصویر لحظه‌ی خطا) و نتیجه‌ی هر تست در فایل‌های نتیجه‌ی
\lr{Allure}
ثبت و سپس به گزارش تعاملی تبدیل می‌شود. این گزارش نقش مستندات اجرایی چارچوب را نیز ایفا می‌کند.

\section{انتگرال‌گیری و تحویل مستمر}
\lr{\gls{CI}}
/ \lr{\gls{CD}}
مجموعه‌ی رویه‌هایی است که در آن‌ها هر تغییر کد به صورت خودکار ساخته، تست و گزارش می‌شود
\cite{crispin2009}. در فصل
\autoref{chap:evaluation}
شکل پایپلاین‌ای که چارچوب را به این چرخه منتقل می‌کند طراحی شده است؛ چارچوب از همان ابتدا با پیش‌نیازهای آن (پیش‌پرواز ساختمان برنامه، برپایی خودکار سرور درایور و خروجی‌های قابل‌انتشار) طراحی شده است.